
\section{Nonasymptotic Analysis}\label{sec:nonasymp}

In this section, we provide the convergence results for both the analytic and numerical schemes in the nonasymptotic regime. In Section~\ref{sec:assumptions}, we lay out our assumptions about our target measure. In Section~\ref{sec:proofstrat}, we introduce the proof strategy for the reader's convenience, which may be useful for extending and proving similar results. In Section~\ref{sec:conv}, we provide the full nonasymptotic bounds -- in particular, in Section~\ref{sec:conv_stationary}, we identify the stationary measure for KIPLD and show an exponential convergence rate to it, as well as, its concentration onto the \gls*{mmle} solution $\bar{\theta}_\star$ as $N$ grows. Following this, in Section~\ref{sec:numerical}, the error bounds for the two algorithms are provided, allowing us to identify the convergence rate of \gls*{kiplmc1} and \gls*{kiplmc2} to the \gls*{mmle} solution $\bar{\theta}_\star$.

\subsection{Assumptions}\label{sec:assumptions}
Our assumptions are generic for the analysis of Langevin diffusions, i.e., we assume strong convexity and $L$-Lipschitz gradients for the potential $U$. These assumptions are akin to the assumptions made in \citep{durmus2017nonasymptotic,dalalyan2017theoretical,dalalyan2019user} for proving the convergence of \gls*{ula}. It is, however, possible to relax them, as can be seen in \citep{zhang2023nonasymptotic,akyildiz2024nonasymptotic}.

We start with the following assumption on the potential $U$.
\begin{assumption}\label{assump:strongconvexity}
Suppose $U\in C^2(\mathbb{R}^{d_\theta+d_x})$. There exists a constant $\mu>0$ such that, for all $z,z'\in \mathbb{R}^{d_\theta+d_x}$,
\begin{equation*}
\langle z-z', \nabla U(z) - \nabla U(z')\rangle \geq \mu \|z-z'\|^2.
\end{equation*}
Since $U\in C^2$, this is equivalent to $\nabla^2 U(z)\succeq \mu I$ for all $z\in \mathbb{R}^{d_\theta+d_x}$, i.e., $U$ is $\mu$-strongly convex.
\end{assumption}

This assumption is equivalent to an assumption of joint strong convexity in $\theta$ and $x$ for $U$, which ensures exponentially fast convergence of the KIPLD to the global minimiser and places a quadratic lower growth bound on $U$.
\begin{assumption}
\label{assump:lipschitz}
There exists a constant $L>0$ such that, for all $z,z'\in \mathbb{R}^{d_\theta+d_x}$,
\begin{equation*}
\|\nabla U(z) - \nabla U(z')\| \leq L \|z-z'\|.
\end{equation*}
\end{assumption}
Together with Assumption~\ref{assump:strongconvexity}, this implies $\mu\leq L$. This assumption is a common assumption in the analysis of stochastic differential equations, as it ensures stability of the diffusion KIPLD, leading to strong solutions, and stable numerical discretisations \cite{kloeden2012numerical}. Observe that this assumption places an upper quadratic growth bound on $U$ or, from the perspective of the probabilistic model, $\log p_\theta(x,y)$ is of quadratic growth in both $\theta$ and $x$. \ref{assump:strongconvexity} and \ref{assump:lipschitz} together ensure that we can apply the Leibniz differentiation rule under integration for the marginal of the probability model $p_\theta(x, y)= e^{-U(\theta,x)}$ \citep[Theorem~16.8]{billingsley1995probability}.

We remark nonetheless that, in the works cited above, the assumptions are on the strong log-concavity of the target measure. In our setting, this corresponds to imposing a structure on the joint statistical model $p_\theta(x, y)$ in $\theta$ and $x$, as our potential is given as $U(\theta,x) = -\log p_\theta(x,y)$. This requires some care on the user side as a statistical model with strongly log-concave posterior $p_\theta(x | y)$ (with fixed $\theta$) may not necessarily satisfy strong log-concavity of $p_\theta(x, y)$ in $(x, \theta)$ jointly. We verify in the experimental section that for some generic models such as Bayesian logistic regression, strict convexity can be shown (which can be improved to strong convexity with regularisation). As noted above, however, these strong assumptions on the statistical model should not be difficult to relax using the standard non-log-concave sampling results -- but this direction is out of scope of the present work which considers the strongly convex case.

\subsection{The proof strategy}\label{sec:proofstrat}
Let $(\theta_n)_{n\geq 0}$ be the sequence of iterates generated by a numerical scheme, for example, \gls*{kiplmc1} or \gls*{kiplmc2}. Let $\pi_\Theta$ denote the $\theta$-marginal of the stationary measure of the diffusion KIPLD, and $\delta_{\bar{\theta}_\star}$ denote the Dirac measure at $\bar{\theta}_\star$. Finally, we denote the law of $\theta_n$ by $\mathcal{L}(\theta_n)$.

In order to proceed, we first note that the optimisation error is given as $\bE[\|\theta_n - \bar{\theta}_\star\|^2]^{1/2} = W_2(\mathcal{L}(\theta_n), \delta_{\bar{\theta}_\star})$. This is due to the fact that the set of couplings between a measure $\nu$ and another measure $\delta_y$ contains a single element $\delta_y \otimes \nu$ \citep[Section~1.4]{santambrogio2015optimal}, thus the infimum in Wasserstein-2 distance is attained by the coupling $\delta_{\bar{\theta}_\star} \otimes \mathcal{L}(\theta_n)$. This implies that $\bE[\|\theta_n - \bar{\theta}_\star\|^2]^{1/2} = W_2(\mathcal{L}(\theta_n), \delta_{\bar{\theta}_\star})$ and consequently using triangle inequality
\small
\begin{align}
\bE[\|\theta_n - \bar{\theta}_\star\|^2]^{1/2} \leq \underbrace{W_2(\mathcal{L}(\theta_n), \pi_\Theta)}_{\textnormal{convergence}}  + \underbrace{W_2(\pi_\Theta, \delta_{\bar{\theta}_\star})}_{\textnormal{concentration}},\label{eq:triangle}
\end{align}
\normalsize
as the Wasserstein distance is a metric. The first term is the convergence of the numerical scheme to the stationary measure $\pi_\Theta$ which will be proved for each scheme separately. The second term in the right-hand side of \eqref{eq:triangle} is the concentration of the stationary measure $\pi_\Theta$ onto the \gls*{mmle} solution $\bar{\theta}_\star$, which is given by Proposition~\ref{prop:mmleerror} below.

\subsection{Convergence bounds}\label{sec:conv}
As outlined above in \eqref{eq:triangle}, we will first show the concentration of the stationary measure $\pi_\Theta$ onto the \gls*{mmle} solution $\bar{\theta}_\star$, then show the convergence of the KIPLD to the stationary measure $\pi_\Theta$. Finally, we will provide the error bounds for the numerical schemes.

\subsubsection{Concentration of the stationary measure}\label{sec:conv_stationary}
This is a non-standard underdamped Langevin diffusion, thus, we need to first identify the stationary measure of the system. Recall that we are only interested in $\theta$-marginal of this stationary measure and we denote it by $\pi_\Theta$. We first have the following proposition.
\begin{proposition} \label{prop:stationary}
Let $\pi_\Theta$ be the $\theta$-marginal of the stationary measure of the KIPLD. Then, we can write its density as
\begin{align}
\pi_\Theta(\theta) \propto \exp( - N \kappa(\theta)).
\end{align}
where $\kappa(\theta) = - \log p_\theta(y)$.
\end{proposition}
\begin{proof} See Appendix~\ref{app:proofstationary}.
\end{proof}
This result shows that the \gls*{kipld} targets the right object: as $N$ grows, $\pi_\Theta$ will concentrate on the minimisers of $\kappa(\theta)$ by a classical result \citep{hwang1980laplace}. This shows that the number of particles $N$ acts as an inverse temperature parameter in the underdamped diffusion. Since the minimiser of $\kappa(\theta)$ is the maximiser of $\log p_\theta(y)$, we can see that the stationary measure of the \gls*{kipld} is concentrating on the \gls*{mmle} solution $\bar{\theta}_\star$. In particular, under the assumption \ref{assump:strongconvexity}, we have the following nonasymptotic concentration result.
\begin{proposition}\label{prop:mmleerror}
Under \ref{assump:strongconvexity}, the function $\theta \mapsto p_\theta(y)$ is $\mu$-strongly log-concave. Furthermore, we have
\begin{equation*}
W_2(\pi_\Theta, \delta_{\bar{\theta}_\star})\leq \sqrt{\frac{d_\theta}{\mu N}},
\end{equation*}
where $\bar{\theta}_\star = \argmax_{\theta\in\mathbb{R}^{d_\theta}} \log p_\theta(y)$,
which is unique.
\end{proposition}

\begin{proof}
See Appendix~\ref{app:proofmmleerror}.
\end{proof}

We note here that this bound is tight, as it is achieved when $\pi_\Theta$ is Gaussian.

\subsubsection{Convergence of the KIPLD to the stationary measure} \label{sec:numerical}
We next demonstrate that the KIPLD converges to its stationary measure $\pi_\Theta$ exponentially fast in the strongly log-concave case.

\begin{proposition}\label{prop:kipld_conv}
Let $(\bm{\theta}_t)_{t\geq 0}$ be the $\theta$-marginal of a solution to the KIPLD, initialised so that
\begin{align*}
\Z_0
:=(\bm\theta_0,N^{-1/2}\X_0^1,\ldots,N^{-1/2}\X_0^N)^\intercal
\sim\nu,
\end{align*}
where $\nu\in\mathcal P(\mathbb R^{d_z})$ has finite second moment and $d_z=d_\theta+Nd_x$.  Independently of $\Z_0$, initialise the mutually independent velocities as
\begin{align*}
\V_0^\theta&\sim\mathcal N(0,N^{-1}\mathbb I_{d_\theta}),&
\V_0^{x_i}\sim\mathcal N(0,\mathbb I_{d_x}),\qquad i\in[N].
\end{align*}
Then, under \ref{assump:strongconvexity} and \ref{assump:lipschitz} and for $\gamma \geq \sqrt{\mu + L}$, for every $t\geq0$ we have
\begin{align*}
W_2(\mathcal{L}(\bm{\theta}_t), \pi_\Theta) \leq \sqrt{2} \exp\left(-\frac{\mu}{\gamma}t\right) \mathbb{E}[\|\Z_0 - \bar{Z}_\star\|^2]^{1/2},
\end{align*}
where $\bar{Z}_\star\sim\widetilde\pi_Z$, the position marginal of the invariant measure $\widetilde\pi$ in Lemma~\ref{lem:scaledstat}, is drawn independently of $\Z_0$.
\end{proposition}
\begin{proof} See Appendix~\ref{app:proofkipld_conv}. \end{proof}

\subsubsection{Nonasymptotic analysis of KIPLMC1}\label{sec:numericalKPLMC1}
In this section we present our first main result, the convergence rate of \gls*{kiplmc1}. 
\begin{theorem}\label{thm:KPLMC1} Let $(\theta_n)_{n\in\bN}$ be the iterates of \gls*{kiplmc1} and suppose that the process is initialised as $Z_0\sim \nu$, where $\nu\in\mathcal{P}(\mathbb{R}^{d_z})$ has bounded second moments, $V^\theta_0\sim \mathcal{N}(0, N^{-1}\mathbb{I}_{d_\theta})$ and $V^{x_i}_0 \sim \mathcal{N}(0, \mathbb{I}_{d_x})$, for $i=1,\dots,N$. Under Assumptions \ref{assump:strongconvexity} and \ref{assump:lipschitz} and with $\gamma\geq \sqrt{\mu+L}$ and $\eta \leq \mu/(4\gamma L)$, we have
\begin{align*}
\bE[\|\theta_n - \bar{\theta}_\star\|^2]^{1/2} &\leq C_1 \left(1-\frac{0.75\mu\eta}{\gamma}\right)^n + \eta C_2 + \sqrt{\frac{C_3}{N}},
\end{align*}
where
\begin{align*}
C_1 &= \sqrt{2} \bE[\|Z_0 - \bar{Z}_\star\|^2]^{1/2},
\qquad C_2 = \sqrt{2}\frac{L}{\mu} \sqrt{\frac{d_\theta + N d_x}{N}}, \qquad
C_3 = \frac{d_\theta}{\mu},
\end{align*}
where $Z_0=(\theta_0, N^{-1/2} X^1_0,\dots, N^{-1/2} X^N_0)^\intercal$, the initialisation step of the \gls*{kiplmc1} where $\bar{Z}_\star\sim\widetilde\pi_Z$ and $\widetilde\pi_Z$ is the position marginal of the invariant measure $\widetilde\pi$ described in Lemma~\ref{lem:scaledstat}.
\end{theorem}
\begin{proof}
See Appendix~\ref{app:proofnumericalKPLMC1}.
\end{proof}

To provide the algorithmic complexity of \gls*{kiplmc1}, we can inspect the bound provided in Theorem~\ref{thm:KPLMC1}. To obtain a complexity result, we first set $N = \mathcal{O}(d_\theta \varepsilon^{-2})$ which ensures the last term in the bound in Theorem~\ref{thm:KPLMC1} is $\mathcal{O}(\varepsilon)$. Next, set $\eta = \mathcal{O}(\varepsilon d_x^{-1/2})$ which ensures that the second term is $\mathcal{O}(\varepsilon)$. It is then easy to see that $n = \widetilde{\mathcal{O}}(\varepsilon^{-1} d_x^{1/2})$ steps are enough to obtain $\bE[\|\theta_n - \bar{\theta}_\star\|^2]^{1/2} \leq \varepsilon$.

\begin{remark}\label{rem:kiplmc1_complexity} This dependency needs to be compared with \gls*{ipla} and \gls*{pgd}. For the former, we note that our dependence to dimension for $N$ is identical to \gls*{ipla} and the dependence of the number of steps is significantly improved, i.e., we require $n = \widetilde{\mathcal{O}}({d_x}^{1/2} \varepsilon^{-1})$ whereas \gls*{ipla} requires $n = \widetilde{\mathcal{O}}({d_x} \varepsilon^{-2})$ steps for the same accuracy. Regarding the latter method, \gls*{pgd}, our dependence on the number of particles is different, as \gls*{pgd} requires $N = \mathcal{O}(d_x \varepsilon^{-2})$ (but this difference is a more general difference stemming from the different analysis techniques, as \gls*{ipla} also has the same dependence as us to the number of particles $N$, see, e.g., Corollary~8 and the subsequent discussion in \citep{caprio2024error}). However, in terms of the number of steps necessary to attain $\varepsilon$-error, we have the same improvement in dimension dependence in $d_x$ and $\varepsilon$ as we have compared to \gls*{ipla}.
\end{remark}

\subsubsection{Nonasymptotic analysis of KIPLMC2}
We present our second main result, showing the convergence rate of \gls*{kiplmc2}.

\begin{theorem}\label{thm:KIPLMC2}

Let $(\theta_n)_{n\in\mathbb{N}}$ be the iterates of \gls*{kiplmc2}, set $d_z=d_\theta+Nd_x$, and define the rescaled initial position
\begin{align*}
Z_0=(\theta_0,N^{-1/2}X_0^1,\ldots,N^{-1/2}X_0^N)^\intercal.
\end{align*}
Suppose that $Z_0\sim\nu$, where $\nu\in\mathcal{P}(\mathbb{R}^{d_z})$ has finite second moment, and, independently of $Z_0$, initialise the mutually independent velocities as
\begin{align*}
V_0^\theta&\sim\mathcal{N}(0,N^{-1}\mathbb{I}_{d_\theta}), &
V_0^{x_i}\sim\mathcal{N}(0,\mathbb{I}_{d_x}),\qquad i\in[N].
\end{align*}
Under Assumptions~\ref{assump:strongconvexity} and~\ref{assump:lipschitz}, if $\gamma\geq2\sqrt{L}$ and $0<\eta\leq\mu/(33\gamma^3)$, then, for every $n\in\mathbb{N}$,
letting $\widetilde\pi_Z$ denote the position marginal of the invariant measure
$\widetilde\pi$ in Lemma~\ref{lem:scaledstat}, we have
\begin{align*}
\bE[\|\theta_n-\bar{\theta}_\star\|^2]^{1/2}
&\lesssim (\sqrt N\,W_2(\nu,\widetilde\pi_Z) +\eta N\sqrt{d_\theta+Nd_x})
\exp\left(-\frac{\mu\eta n}{6\gamma}\right)\\
&\quad+\eta\sqrt{N(d_\theta+Nd_x)}
+\sqrt{\frac{d_\theta}{N}}.
\end{align*}
The implicit constant depends only on $\mu$, $L$, and $\gamma$; in particular,
it is independent of $n$, $N$, $\eta$, $d_\theta$, and $d_x$.
\end{theorem}

\begin{proof}
See Appendix \ref{app:proofnumericalKPLMC2}.
\end{proof}


To obtain complexity, similarly to the analysis provided for \gls*{kiplmc1}, we first set $N\asymp d_\theta\varepsilon^{-2}$ to obtain an $\mathcal{O}(\varepsilon)$ bound for the last term in Theorem~\ref{thm:KIPLMC2}.  With this choice,
\begin{align*}
\sqrt{N(d_\theta+Nd_x)}
\asymp\frac{d_\theta}{\varepsilon}
\sqrt{1+\frac{d_x}{\varepsilon^2}}.
\end{align*}
We therefore choose
\begin{align*}
\eta
&\asymp\min\left\{\frac{\mu}{33\gamma^3},
\frac{\varepsilon}{\sqrt{N(d_\theta+Nd_x)}}\right\}\\
&\asymp\min\left\{\frac{\mu}{33\gamma^3},
\frac{\varepsilon^2}
{d_\theta\sqrt{1+d_x/\varepsilon^2}}\right\},
\end{align*}
which makes the nondecaying discretisation term of order $\varepsilon$.  In the small-$\varepsilon$ regime, this gives $\eta=\mathcal{O}(\varepsilon^3/(d_\theta\sqrt{d_x}))$.  Finally, writing $D_0:=W_2(\nu,\widetilde\pi_Z)$ and taking
\begin{align*}
n\gtrsim\frac{1}{\eta}
\log\left(1+\frac{\sqrt N\,D_0
+\eta N\sqrt{d_\theta+Nd_x}}{\varepsilon}\right)
\end{align*}
controls both exponentially decaying terms.  Hence, assuming that the initial Wasserstein distance grows at most polynomially in $N$ and the dimensions, $n=\widetilde{\mathcal{O}}(d_\theta\sqrt{d_x}\,\varepsilon^{-3})$ steps are sufficient to obtain an accuracy of order $\varepsilon$, where the notation also suppresses factors depending only on $\mu$, $L$, and $\gamma$.

\begin{table}[t]
    \centering
    \def \arraystretch{1}
    \resizebox{.5\linewidth}{!}{\begin{tabular}{c|c c c}
    \hline
         &  step count ($n$) & $N$ & $\eta$\\
        \hline
       KIPLMC1  &  $\widetilde{\mathcal{O}}(\sqrt{d_x}/\varepsilon)$ & $\mathcal{O}(d_\theta/\varepsilon^2)$ & $\mathcal{O}(\varepsilon/\sqrt{d_x})$ \\
       KIPLMC2  &  $\widetilde{\mathcal{O}}(d_\theta\sqrt{d_x}/\varepsilon^3)$ & $\mathcal{O}(d_\theta/\varepsilon^2)$ & $\mathcal{O}(\varepsilon^3/(d_\theta\sqrt{d_x}))$\\
       IPLA & $\widetilde{\mathcal{O}}(d_x/ \varepsilon^{2})$ & $\mathcal{O}(d_\theta/\varepsilon^2)$ & $\mathcal{O}(\varepsilon^2/d_x)$\\
       PGD & $\widetilde{\mathcal{O}}(d_x/\varepsilon^2)$ & $\mathcal{O}(d_x/\varepsilon^2)$ & $\mathcal{O}(\varepsilon^2 /d_x)$\\
       \hline
    \end{tabular}}
    \caption{Comparison of orders required to achieve an error of order $\varepsilon$ in $W_2$, for step count, particle number $N$ and step size $\eta$.}
    \label{tab:err_rate}
\end{table}

A few comments are in order for this complexity analysis for \gls*{kiplmc2}.
\begin{remark}\label{rem:kiplmc2_complexity} In general, to avoid poor dependence on $N$, the error bounds of numerical discretisations must depend on $L$ and $\mu$ with the matching powers. This ensures that the factors cancel out appropriately, preventing any undesirable effects. This is achieved by the analysis of the exponential integrator in \citep{dalalyan2018sampling} with a linear dependence on the condition number $L/\mu$, thus our \gls*{kiplmc1} bound does not have a bad dependence to $N$. This enables us to obtain improved algorithmic complexity as discussed in Remark~\ref{rem:kiplmc1_complexity}. However, for \gls*{kiplmc2}, the OBABO scheme does not have this dependence on the condition number (see, e.g., \citep[Table~2]{monmarche2021high}) under just \ref{assump:strongconvexity} and \ref{assump:lipschitz}. This results in a bound in Theorem~\ref{thm:KIPLMC2} that grows as $N$ grows, however, we note that it is the authors' belief that this dependence may be an artifact of the proof method employed, rather than reflecting a deteriorating accuracy with growing $N$ for \gls*{kiplmc2}. This is observed in the numerical experiments we present where increasing $N$ does not lead to a growing error.
\end{remark}


See Table \ref{tab:err_rate} for a summary of our discussion.
